% Partie : intuitionistic-logic
% Chapitre : tableaux
% Section : rules

\documentclass[../../../include/open-logic-section]{subfiles}

\begin{document}

\olfileid[fr]{int}{tab}{rul}

\olsection{Règles pour la logique intuitionniste}

Les règles pour les connecteurs $\land$ et $\lor$ sont les mêmes que
dans les !!{tableau}s propositionnels signés ordinaires, avec des
préfixes en plus. Dans chaque cas, la règle appliquée à une
!!{formula} signée $\sFmla{S}{!A}[\sigma]$ produit de nouvelles
!!{formula}s portant elles aussi le préfixe~$\sigma$. Cela se
comprend intuitivement : si $!A \land !B$ est vraie au monde désigné
par~$\sigma$, alors $!A$ et $!B$ sont vraies au même monde~$\sigma$. Les
règles pour $\land$ et $\lor$ figurent dans le
\olref{tab:prop-rules}.

\begin{table}
  \[\def\arraystretch{3}\begin{array}{|c|c|}
    \hline
    \AxiomC{\sFmla{\True}{!A \land !B}[\sigma]}
    \RightLabel{\TRule{\True}{\land}}
    \UnaryInfC{\sFmla{\True}{!A}[\sigma]}
    \noLine
    \UnaryInfC{\sFmla{\True}{!B}[\sigma]}
    \DisplayProof
    &
    \AxiomC{\sFmla{\False}{!A \land !B}[\sigma]}
    \RightLabel{\TRule{\False}{\land}}
    \UnaryInfC{$\sFmla{\False}{!A}[\sigma] \quad \mid \quad
      \sFmla{\False}{!B}[\sigma]$}
    \DisplayProof
    \\[2ex]
    \hline
    \AxiomC{\sFmla{\True}{!A \lor !B}[\sigma]}
    \RightLabel{\TRule{\True}{\lor}}
    \UnaryInfC{$\sFmla{\True}{!A}[\sigma] \quad \mid \quad
      \sFmla{\True}{!B}[\sigma]$}
    \DisplayProof
    &
    \AxiomC{\sFmla{\False}{!A \lor !B}[\sigma]}
    \RightLabel{\TRule{\False}{\lor}}
    \UnaryInfC{\sFmla{\False}{!A}[\sigma]}
    \noLine
    \UnaryInfC{\sFmla{\False}{!B}[\sigma]}
    \DisplayProof
    \\[2ex]
    \hline
  \end{array}\]
  \caption{Règles des !!{tableau}s préfixés pour $\land$ et $\lor$}
  \ollabel{tab:prop-rules}
\end{table}

La condition de fermeture ressemble à celle des !!{tableau}s
ordinaires, mais les préfixes, et pas seulement les !!{formula}s,
doivent aussi être pris en compte. On
peut même être un peu moins restrictif : dans un modèle
intuitionniste, une !!{formula} vraie reste vraie à tout monde
accessible. Il est donc impossible que $!A$ soit vraie en~$\sigma$
et fausse en un préfixe accessible~$\sigma.{*}$. Une branche est
ainsi fermée si elle contient à la fois
\[
\sFmla{\True}{!A}[\sigma] \quad\text{et}\quad \sFmla{\False}{!A}[\sigma.{*}]
\]
pour un préfixe~$\sigma$ et une !!{formula}~$!A$. Si les signes
sont inversés, c'est-à-dire si la branche contient
\[
\sFmla{\False}{!A}[\sigma] \quad\text{et}\quad \sFmla{\True}{!A}[\sigma.{*}]
\]
cette paire ne ferme la branche que si $*$ désigne la suite vide.

Une branche est également fermée si elle contient
$\sFmla{\True}{\bot}[\sigma]$.

Les hypothèses sont introduites comme dans les !!{tableau}s
ordinaires, mais reçoivent toujours le préfixe~$1$. (Le choix du
préfixe importe peu, pourvu qu'il soit le même pour toutes les
hypothèses.) Ainsi, nous disons que
\[
!B_1, \dots, !B_n \Proves !A
\]
si et seulement s'il existe un tableau fermé pour les hypothèses
\[
\sFmla{\True}{!B_1}[1], \dots, \sFmla{\True}{!B_n}[1],
\sFmla{\False}{!A}[1].
\]

Pour le conditionnel~$\lif$, les règles diffèrent des règles
classiques et modales. La règle $\TRule{\True}{\lif}$ prolonge une
branche contenant $\sFmla{\True}{!A \lif !B}[\sigma]$ par
$\sFmla{\False}{!A}[\sigma.{*}]$ sur une branche et
$\sFmla{\True}{!B}[\sigma.{*}]$ sur l'autre. Elle ne peut être
appliquée qu'à un préfixe~$\sigma.{*}$ qui figure \emph{déjà} sur la
branche à laquelle elle s'applique. Appelons un tel préfixe
« utilisé » sur cette branche. (Comme $\sigma.{*}$ comprend
$\sigma$ lui-même, on peut toujours appliquer la règle en ajoutant
$\sFmla{\False}{!A}[\sigma]$ et
$\sFmla{\True}{!B}[\sigma]$ sur deux branches distinctes.)

La règle $\TRule{\False}{\lif}$ prolonge une branche contenant
$\sFmla{\False}{!A \lif !B}[\sigma]$ par les deux formules
$\sFmla{\True}{!A}[\sigma.n]$ et
$\sFmla{\False}{!B}[\sigma.n]$ sur la même branche, où
$\sigma.n$ est un préfixe nouveau pour celle-ci.

Les règles pour $\lnot$ sont définies de manière analogue, en
utilisant la définition de $\lnot !A$ comme $!A \lif \lfalse$.

Ces règles figurent dans le \olref{tab:rules-lif-lnot}.

\begin{table}
  \[\def\arraystretch{3}\begin{array}{|c|c|}
    \hline
    \AxiomC{\sFmla{\True}{\lnot !A}[\sigma]}
    \RightLabel{\TRule{\True}{\lnot}}
    \UnaryInfC{$\sFmla{\False}{!A}[\sigma.{*}]$}
    \DisplayProof
    &
    \AxiomC{\sFmla{\False}{\lnot !A}[\sigma]}
    \RightLabel{\TRule{\False}{\lnot}}
    \UnaryInfC{\sFmla{\True}{!A}[\sigma.n]}
    \DisplayProof
    \\[1ex]
    \text{$\sigma.{*}$ est utilisé} & \text{$\sigma.n$ est nouveau}\\
    \hline
    \AxiomC{\sFmla{\True}{!A \lif !B}[\sigma]}
    \RightLabel{\TRule{\True}{\lif}}
    \UnaryInfC{$\sFmla{\False}{!A}[\sigma.{*}] \quad \mid \quad
      \sFmla{\True}{!B}[\sigma.{*}]$}
    \DisplayProof
    &
    \AxiomC{\sFmla{\False}{!A \lif !B}[\sigma]}
    \RightLabel{\TRule{\False}{\lif}}
    \UnaryInfC{\sFmla{\True}{!A}[\sigma.n]}
    \noLine
    \UnaryInfC{\sFmla{\False}{!B}[\sigma.n]}
    \DisplayProof
    \\[1ex]
    \text{$\sigma.{*}$ est utilisé} & \text{$\sigma.n$ est nouveau}\\
    \hline
  \end{array}\]
  \caption{Règles des !!{tableau}s préfixés pour $\lnot$ et $\lif$}
  \ollabel{tab:rules-lif-lnot}
\end{table}

\end{document}
